% Source: https://artofproblemsolving.com/community/c3h2315509_3_year_anniversary
% Stable identity: aops-topic-2315509-q06
Take the cross-section through the axis of the cone. Let $V$ be the vertex and $\theta$ the half-angle of the cone. A ball of radius $r$ tangent to the sides has its centre on the axis at distance $d=\dfrac{r}{\sin\theta}$ from $V$ (the radius to the point of tangency is perpendicular to the side).

\begin{asy}
usepackage("amsmath");
size(4cm);
defaultpen(fontsize(10pt));
real s=1/3, c=sqrt(8)/3; real t=s/c;   // half-angle theta with sin=1/3
pair O=(0,0);
real H=8;
draw(O--(H*t,H)); draw(O--(-H*t,H)); draw((-H*t,H)--(H*t,H),linewidth(0.9));
real[] r={0.5,1,2};
for(int k=0;k<3;++k){ real d=3*r[k]; draw(circle((0,d),r[k]),blue); dot((0,d),blue); }
draw((0,0)--(0,8),dashed+gray);
pair P2=(0,6), T2=(0,6)+2*dir(180+asin(1/3)*180/pi);
draw(P2--T2,red); label("$2$",(P2+T2)/2,N,red+fontsize(8pt));
draw((3.3,0)--(3.3,6),Arrows(3)); label("$6$",(3.3,3),E,fontsize(8pt)); draw((3.3,6)--(3.3,8),Arrows(3)); label("$2$",(3.3,7),E,fontsize(8pt));
draw((0,6)--(3.4,6),dotted); draw((0,0)--(3.4,0),dotted);
draw(arc(O,0.75,90,90+asin(1/3)*180/pi)); label("$\theta$",O+0.8*dir(90+asin(1/3)*180/pi),W,fontsize(8pt));
label("$V$",O,S); label("rim",(-H*t,H),W);
\end{asy}

\textbf{The angle.} For consecutive scoops of radii $r$ and $2r$ that touch, the distance between the centres is $r+2r=3r$, and it is also $\dfrac{2r}{\sin\theta}-\dfrac{r}{\sin\theta}=\dfrac{r}{\sin\theta}$. So $\sin\theta=\dfrac13$.

\textbf{The height.} The largest scoop has radius $2$, so its centre is $\dfrac{2}{\sin\theta}=6$ inches above $V$, and its top is $6+2=8$ inches above $V$. (It touches the sides at height $6\cos^2\theta=\dfrac{16}{3}<8$, so inside the cone.) Since the top of this scoop is level with the rim, the height of the cone is $\boxed{8\text{ inches}}$.
